\scanpage{70}
% archive sheet number: 69
\uline{Exemples}

\textcircled{1} $d'=d''=\varnothing$. \deleted{Alors} a) Prenant $\Delta=(\mathrm{Cat}_d)$, les Pbs 1) et 2) ont une solution évidente resp. bien connue. Les théories Algébriques sont exactement celles qui sont représentées par des \uline{petites} catégories. Les $I\longrightarrow R$ générateurs sont uniquement les ess. surjectifs.

b) \uline{Condition} d'exactitude. On peut par exemple prendre \uline{$\Delta=(\text{groupoïdes})$}.
\margin{\diagram{$\underline{\unsure{\rho|p}}=\left[\,\begin{tikzcd}[ampersand replacement=\&,column sep=small]
\overset{a}{\bullet}\arrow[r,"u"] \& \overset{b}{\bullet}
\end{tikzcd}\,\right]$}
\unsure{$u$ rendu un|on rend $u$ un} isom}
Les \deleted{\unsure{cas}} Pb 1 a une solution connue (groupoïde fondamental d'une catégorie), le Pb 2) est trivial. Les théories Alg. sont celles repr. par des petites \deleted{catégories} \inserted{groupoïdes}. Les $I\longrightarrow R$ générateurs sont les ess. surj.

c) Prenons $\Delta=(\text{catégories ordonnées})$. Pb 1 a une solution évidente (\deleted{\unsure{catégori|catégorie}} ens. préordonné associé à une catégorie), Pb 2 une solution facile.
\margin{\diagram{$\underline{\unsure{\rho|p}}=\left[\,\begin{tikzcd}[ampersand replacement=\&,column sep=small]
\underset{a}{\bullet}\arrow[r,shift left,"u"]\arrow[r,shift right,"v"'] \& \underset{b}{\bullet}
\end{tikzcd}\,\right]$}
\deleted{\unsure{pour|que}} \unsure{ici|là} on rend $u$ égal à $v$}
Les théories algébriques sont \inserted{celles} repr. par petits ens. \unsure{(pré)ordonnés|(pré), ordonnés}. Les $I\longrightarrow R$ générateurs sont ess. surjectifs.
